\chapter{Errata e incongruenze nelle fonti}\label{errata-cap}

\begin{nota}[A cosa serve]
Piccoli errori delle soluzioni e delle slide del professore, trovati confrontando i testi con i ricalcoli. In sede d'esame conta il \textbf{metodo}; ma conoscere l'errore evita di ``correggere'' una formula giusta o di copiarne una sbagliata.
\end{nota}

%==========================================================
\section{Soluzioni dei temi}
\begingroup\small
\renewcommand{\arraystretch}{1.3}
\begin{tabularx}{\linewidth}{>{\raggedright\arraybackslash}p{2.6cm}X}
\toprule
Tema & Incongruenza e correzione\\
\midrule
2025 -- simulazione (9 CFU) & La soluzione stampata schedula $B\to A$ con media 0.5 s e $A\to B$ con media 1 s: \textbf{scambiate}. Giusto: dopo $A\to B$ si schedula $B\to A$ con media $1/\gamma_B=1$ s; dopo $B\to A$ si schedula $A\to B$ con media $1/\gamma_A=0.5$ s.\\
2020 -- Es.~1 & Si cita ``sostituendo nella 3a equazione'' ma è la normalizzazione; rimando a un ``punto 8'' che non esiste. Risultati corretti.\\
2020 -- Es.~3 & Refuso ``tra 43 e $-2$'': leggere $-4$ e $-2$.\\
Gen.\ variabili casuali, foglio 2 c) & Geometrica di media 8 con $X=\lfloor\ln U/\ln(1-p)\rfloor$, $p=\tfrac18$: questa variabile parte da 0 e ha media $(1-p)/p=7$. Per supporto $\{1,2,\dots\}$ e media 8: $\lceil\ln U/\ln(1-p)\rceil$, ovvero floor$+1$. Controlla la convenzione del testo.\\
\bottomrule
\end{tabularx}

\begin{tabularx}{\linewidth}{>{\raggedright\arraybackslash}p{2.6cm}X}
\toprule
Gen.\ variabili casuali, foglio 1 c) & L'integrale ha estremo inferiore $-1$ stampato per errore (deve essere 0). Risultato $X=\sqrt[3]{3U}$ corretto.\\
Gen.\ variabili casuali, foglio 3 d) & Il testo dice ``discreta'', la soluzione dà la formula continua $-5\ln U$.\\
Sistema Completo 2 & Chiama $1/x^2$ ``Pareto con shape 2'': in convenzione standard lo shape è $\alpha=1$ (media infinita).\\
2005-06-01 -- Es.~1 & Non è data la lunghezza media dei pacchetti: si lavora in kbit/s e si dichiara l'assunzione. La soluzione non è nella cartella (ricalcolata\mio).\\
2005-06-14 -- Es.~2 & Usa $-\ln U$ invece di $-\ln(1-U)$: giustificato, $U$ e $1-U$ sono equidistribuiti.\\
\bottomrule
\end{tabularx}

\begin{tabularx}{\linewidth}{>{\raggedright\arraybackslash}p{2.6cm}X}
\toprule
2006-05-31 -- Es.~2 & Il testo scrive ``4 failures'' ma i calcoli usano giustamente 3.75/anno e MTTF$=2336$ h.\\
2006-05-31 -- Es.~3 c) & ``media pari a $-\mu$'' è un refuso: media $\mu$, $X=-\mu\ln U$. Soluzione assente (ricalcolata\mio).\\
2006-06-13 -- Es.~1 & ``$\lambda_a=\lambda-\lambda_a$'' va letto $\lambda-\lambda_r$; il risultato $\lambda_a=\lambda/2$ è giusto.\\
2006-06-27 -- Es.~3 & $A_{2:1}$ riportata $0.99999874$; esatto $0.99999873$ (differenza irrilevante).\\
2007-05-31 -- Es.~3 & $T_{\min}=0.01321$ s dichiarato; con la formula approssimata e $X=2828$ si ottiene $\mathbf{0.01313}$ s. Senza approssimazione ($\lceil\cdot\rceil$ esatto): $X\approx2858$, $T=0.01303$ s.\\
\bottomrule
\end{tabularx}

\begin{tabularx}{\linewidth}{>{\raggedright\arraybackslash}p{2.6cm}X}
\toprule
2007-06-12 -- Es.~2, 3 & Soluzioni assenti: ricalcolate\mio\ ($T=96.2$ ms; con il nuovo flusso C$\to$D: diretto $112.4$ ms, via CBD $140.4$ ms; $A_{1:1}$ non Gold, 2:1 Gold).\\
2008 -- Es.~3 & Il punto 3 riprende numeri del 2005 (``$1-0.9966$'', ``$1.1648\cdot10^{-5}$''). Valori corretti: $U=0.00273$ (Bronze), $U_{1:1}=7.465\cdot10^{-6}$ (Premium).\\
2009 -- Es.~2 & Soluzione assente (ricalcolata\mio): $T=86.4$ ms, $\eta_{\max}=4/3$, $T(\eta\to0)=25.15$ ms.\\
E$_2$|M|1|0 & Il file contiene solo il disegno della catena, non lo svolgimento (completato\mio).\\
Link orientati & Il verso di BC/CB cambia fra i temi (B$\to$C nel 2005--06, C$\to$B nel 2007--09): leggere sempre la figura.\\
\bottomrule
\end{tabularx}

\begin{tabularx}{\linewidth}{>{\raggedright\arraybackslash}p{2.6cm}X}
\toprule
2007-05 (arrivi scoraggiati) & Il blocco è calcolato come $p_C$ (prob.\ dello stato pieno) anche se il tasso di arrivo in quello stato è ridotto: è la convenzione del professore.\\
\bottomrule
\end{tabularx}
\endgroup

%==========================================================
\section{Slide}
\begingroup\small
\renewcommand{\arraystretch}{1.3}
\begin{tabularx}{\linewidth}{>{\raggedright\arraybackslash}p{3.2cm}X}
\toprule
Slide & Incongruenza\\
\midrule
Reti di code, 16--18 (grado di servizio, 6 nodi) & La tabella elenca \textbf{una sola direzione} dei link; essendo i link full-duplex e la matrice simmetrica, la direzione opposta porta lo stesso carico: $\sum\lambda_i=328$ su 16 link orientati, $T^*=\tfrac{14.2}{124}=114$ ms (giusto). Il ``numero medio di link attraversati $=164/124=1.32$'' è invece $328/124\approx2.65$; con soli 8 link si otterrebbe 57 ms, che è il $T_0$ per $\eta\to0$ calcolato correttamente.\\
Congestione, 26 (CQL) & La slide stampa $b_M=\sqrt{B/N}$; la forma usuale è $b_M=B/\sqrt N$. Da verificare sugli appunti.\\
Cellular Networks, 42 (GSM) & Raggio 37.8 km ma il calcolo usa 35 km: $2\tau=233\ \mu$s $\approx63$ bit con 35 km; $252\ \mu$s $\approx68.25$ bit con 37.8 km.\\
Sistemi a coda, 26 (Erlang) & La densità stampata ha media $n/\lambda$, le statistiche sono quelle di un Erlang di media $1/\lambda$ con stadi di tasso $n\lambda$ (slide 29): usare quest'ultima convenzione.\\
Simulazione, 76 (intervalli sulla differenza) & La formula riporta gli estremi $x_l-y_l$, $x_u-y_u$: gli estremi corretti sono $x_l-y_u$ e $x_u-y_l$.\\
Simulazione, 54 & $\lambda$ indica la \textbf{media} dell'esponenziale, nelle slide di code è il \textbf{tasso}.\\
ARQ, 44 & $a=t_T/t_I\ge1$ include $t_I$ stesso: non è il solito $t_{\text{prop}}/t_I$.\\
\bottomrule
\end{tabularx}
\endgroup

%==========================================================
\section{Cosa è nelle fonti e cosa è un completamento}
Nella guida i contenuti non presenti nelle slide o nelle soluzioni del professore sono marcati con \mio. Elenco completo:
\par\smallskip
\begingroup\small
\renewcommand{\arraystretch}{1.3}
\begin{tabularx}{\linewidth}{>{\raggedright\arraybackslash}p{4.6cm}X}
\toprule
Contenuto & Origine\\
\midrule
Rejection (Von Neumann), composizione, generazione Erlang e iperesponenziale & Non nelle slide di Simulazione (che trattano inversione, discreto, esponenziale, Rayleigh, gaussiana, geometrica, Poisson). Aggiunti\mio\ come standard.\\
Serie/parallelo per 1+1 e $U_{1:1}=(1-A)^2$ & Le slide danno solo il modello 1:$N$/M:$N$; la regola serie/parallelo è aggiunta\mio.\\
Funzione potenza M/M/1 ($P=\lambda(\mu-\lambda)$, $\lambda^*=\mu/2$) & Slide: solo la frase ``massima a $N=1$'' e il tema 14/06/2005; il calcolo è dalla soluzione.\\
Catena E$_2$|M|1|0 e $\Pb=\big(\tfrac{2\lambda}{2\lambda+\mu}\big)^2$ & Solo il disegno nel file del professore; svolgimento\mio.\\
Soluzioni ricalcolate: 2005-06-01 Es.~1, 2006-05-31 Es.~3, 2007-06-12 Es.~2 e 3, 2009 Es.~2 & Assenti nella cartella: ricalcolate\mio\ e verificate con \texttt{verifica\_numeri.py}.\\
Derivata ottima di flow deviation per il ritardo totale ($\mu/(\mu-\lambda)^2$) & Le slide usano $\partial T_i/\partial\lambda_i$; la variante è un'osservazione\mio.\\
Erlang C, Engset oltre le slide 60--62, PASTA come nome & Non nelle slide; PASTA è usata implicitamente.\\
TCP (AIMD, goodput), routing (OSPF/RIP), cellulari, OPL/CPLEX & \textbf{Non inclusi}: non compaiono negli esercizi dei temi.\\
\bottomrule
\end{tabularx}
\endgroup
